\chapter{Machine Learning}
In diesem Teil werden die benötigte Datenreinigung, verschiedene Machine-Learning-Algorithmen und schließlich das Ergebnis präsentiert.

\section{Datenreinigung}
Die gemessenen Rohdaten sind in Tabelle \ref{tab:trainingsdaten} dargestellt. 

\bgroup
\def\arraystretch{1.5}
	\begin{table}[h]
		\centering
		\resizebox{\textwidth}{!}{
		\begin{tabular}{|c|c|c|c|c|c|c|c|c|c|c|}
			\hline
			id & acc\_x & acc\_y & acc\_z & gyro\_x & gyro\_y & gyro\_z & mag\_x & mag\_y & mag\_z & timestamp \\
			\hline
			68... & -0.62 & -0.05 & 0.13 & -38.39 & 23346 & 36434 & 91 & 93 & 243 & 2025-05-20 12:22:03.646+0000 \\
			\hline
			ed... & -1.87 & 41275 & 0.21 & 24.23 & -21489999 & -33.69 & 10 & 71 & 212 & 2025-05-20 12:24:44.838+0000 \\
			\hline
			4e... & -0.66 & -0.55 & 0.02 & 40.95 & -13.37 & 0.43 & 75 & 104 & 260 & 2025-05-20 12:22:04.183+0000 \\
			\hline
			22... & -2.65 & -0.47 & -0.1 & -43.21 & 175.17 & 107.79 & -52 & 111 & 300 & 2025-05-20 12:23:17.673+0000 \\
			\hline
			1d... & -1.09 & -0.53 & 0.01 & -60.91 & -20.14 & -26.37 & 57 & 78 & 261 & 2025-05-20 12:20:25.462+0000 \\
			\hline
		\end{tabular}
		}
		\caption{gemessene Daten}
		\label{tab:trainingsdaten}
	\end{table}
\egroup

Bevor diese Daten dem überwachten Machine-Learning-Algorithmus übergeben werden können, müssen sie zunächst aufbereitet und gelabelt werden. Dies umfasst die Bereinigung der Rohdaten, die Aufteilung in geeignete Zeitfenster und die Zuordnung von Labels zu den jeweiligen Aktivitäten. 

\subsection{Labeling}
Labels sind essenziell, da sie dem überwachten Lernalgorithmus die Unterscheidung und Klassifizierung der verschiedenen Aktivitäten ermöglichen. Da die Daten unsortiert in der Datenbank gespeichert werden, muss während der Datenerfassung erfasst werden, zu welchen Zeiten welche Aktivität durchgeführt wurde. Mithilfe der Zeitstempel kann anschließend ermittelt werden, welche Aktivität zu dieser Zeit ausgeführt wurde. Die gewählten Labels sind die Aktivitäten, die am Ende klassifiziert werden sollen, also "Gehen\dq, "Joggen\dq\, und "Rennen\dq. Die Daten werden mithilfe eines Python-Skripts verarbeitet. Dabei werden alle zu einem Durchlauf gehörenden Daten in einer separaten CSV-Datei gespeichert. Dadurch entstehen für jede Aktivität, wie beispielsweise "Gehen”, mehrere CSV-Dateien, die jeweils einen vollständigen Lauf enthalten. Dies entspricht Schritt 1, wie in Abbildung \ref{fig:pakete} zu sehen ist.

\subsection{Aufteilung in 5 Sekunden Pakete}
Das Ziel besteht darin, ein System zu entwickeln, das die aktuelle Bewegung einer Person erkennt. Rennt die Person beispielsweise, soll das System „Rennen” ausgeben. Die Herausforderung besteht darin, das System entsprechend zu trainieren. Zwar werden für eine Aktivität wie "Rennen\dq\, 2-minütige Daten aufgenommen, jedoch soll das System nicht erst nach 2 Minuten eine Antwort herausgeben, sondern sofort. Das System soll die Bewegung der Person möglichst schnell erkennen. Deswegen wurde überlegt, das System nicht auf die vollen zwei Minuten zu trainieren, sondern kleinere Pakete dieser zwei Minuten zu erstellen. Um die Paketgröße zu bestimmen, wurde Folgendes getan: Würde man nur Daten aus einer Sekunde nehmen, gäbe es nur zwei Datenpunkte pro Paket, da der Arduino alle 0,5 Sekunden Daten sendet.  Bei nur zwei Datenpunkten kann das System nur schwer eine Bewegung oder ein Muster erkennen. Deswegen wurden fünf Sekunden gewählt. Bei fünf Sekunden sind es zehn Datenpunkte, was ausreicht, um ein Muster zu erkennen. Zugleich ist es kurz genug, um zu sagen, was die Person im Moment gerade macht. Im Endsystem kann das System also 5 Sekunden lang Daten messen und diese dann dem Machine-Learning-Algorithmus schicken, der sie bearbeitet und dem User die Bewegung mitteilt, die diese gerade ausführt. In Abbildung \ref{fig:pakete} ist durch den roten Strich zu erkennen, wie die ersten fünf Sekunden der ersten gelabelten CSV-Datei entnommen wurden, um daraus ein gelabeltes Paket (Schritt 3) zu generieren.

\begin{figure}[h]
	\centering
	\includegraphics[width=\linewidth]{pakete}
	\caption{Schritte zur Paketaufteilung  (eigene Abbildung)}
	\label{fig:pakete}
\end{figure}

\section{Wahl des Machine Learning Modells}
Für die Klassifizierung der Bewegungsdaten wurden verschiedene Machine Learning Modelle evaluiert. Die Wahl des geeigneten Modells hängt von mehreren Faktoren ab, darunter die Genauigkeit, die Trainingszeit und die Eignung für Echtzeitanwendungen.

\subsection{Evaluierte Modelle}
Im Folgenden werden die meistgenutzten überwachten Machine-Learning-Modelle betrachtet und ihre jeweiligen Vor- und Nachteile analysiert, um das passende Modell zu finden.

\subsubsection{Decision Trees}
Decision Trees sind einfache und interpretierbare Modelle, die gut für Klassifizierungsaufgaben geeignet sind. Sie können jedoch zu Overfitting neigen und sind möglicherweise nicht die beste Wahl für komplexe Datensätze.

\subsubsection{Random Forest}
Random Forest ist ein Ensemble-Lernverfahren, das mehrere Decision Trees kombiniert, um die Genauigkeit zu erhöhen und Overfitting zu reduzieren. Es ist robust und funktioniert gut mit großen Datensätzen.

\subsubsection{Support Vector Machines (SVM)}
SVMs sind leistungsstarke Modelle für die Klassifizierung, insbesondere für hochdimensionale Daten. Sie können jedoch rechenintensiv sein und sind möglicherweise nicht ideal für Echtzeitanwendungen.

\subsubsection{Neural Networks}
Neurale Netze, insbesondere Convolutional Neural Networks (CNNs) und Recurrent Neural Networks (RNNs), sind in der Lage, komplexe Muster in den Daten zu erkennen. Sie erfordern jedoch eine große Menge an Trainingsdaten und Rechenleistung.

\subsection{Gewähltes Modell}
In Tabelle 3 ist noch einmal zu sehen, wie die verschiedenen Modelle hinsichtlich Genauigkeit, Trainingszeit und Echtzeitfähigkeit abschneiden.

\bgroup
\def\arraystretch{1.5}
\begin{table}[h]
	\centering
	\resizebox{\textwidth}{!}{
		\begin{tabular}{|c|c|c|c|}
			\hline
			Modell & Genauigkeit & Trainingszeit & Echtzeitfähigkeit \\
			\hline
			Decision Tree & Mittel & Kurz & Hoch \\
			\hline
			Random Forest & Hoch & Mittel & Hoch \\
			\hline
			SVM & Hoch & Lang & Niedrig \\
			\hline
			Neural Network & Sehr Hoch & Sehr Lang & Niedrig \\
			\hline
		\end{tabular}
	}
	\caption{Vergleich der evaluierten Machine Learning Modelle}
	\label{tab:modellvergleich}
\end{table}
\egroup

Für dieses Projekt wird ein Modell benötigt, das eine hohe Genauigkeit aufweist und nicht zu Overfitting neigt. Decision Trees sind hierfür weniger geeignet, da sie bei begrenzten Trainingsdaten zu Overfitting neigen. Neuronale Netzwerke scheiden ebenfalls aus, da die verfügbare Menge an Trainingsdaten nicht ausreicht. Die Trainingszeit ist weniger kritisch, da das Training nur einmal durchgeführt wird. Entscheidend ist jedoch eine hohe Echtzeitfähigkeit, da alle fünf Sekunden neue Daten verarbeitet werden müssen. Daher eignen sich SVM und neuronale Netzwerke eher weniger. Am besten zu diesen Kriterien passt das Random-Forest-Modell, da es eine gute Balance zwischen Genauigkeit und Rechenaufwand bietet. Es ist robust gegenüber Overfitting und kann gut mit den vorliegenden Bewegungsdaten umgehen.

\section{Implementierung und Ergebnis}
In diesem Abschnitt werden die Implementierung und das Ergebnis des Lernverfahrens präsentiert.

\subsection{Implementierung}
Das Random-Forest-Modell wurde mit der Python-Bibliothek „scikit-learn” implementiert. Die Trainingsdaten werden in Daten, mit denen das Modell trainiert wird, und Daten, mit denen das Modell anschließend getestet wird, um die Leistung des Modells zu evaluieren, unterteilt. Dies erfolgt mit dem folgenden Code:

\begin{lstlisting}[language=Python]
	X_train, X_test, y_train, y_test = train_test_split(X, y_encoded, test_size=0.2, random_state=42)
\end{lstlisting}

\subsection{Ergebnis des Trainings}
Nachdem das Modell trainiert wurde, kann ein Klassifikationsbericht erstellt werden, um dessen Leistung zu bewerten. Um diesen Bericht zu verstehen, müssen zunächst die verschiedenen Begriffe erläutert werden.

\subsubsection{Begriffe zum Klassifikationsbericht}
\begin{itemize}
	\item \textbf{precision (Präzision)}: Der Anteil der korrekt vorhergesagten positiven Instanzen an allen vorhergesagten positiven Instanzen. Eine hohe Präzision bedeutet, dass das Modell wenige falsch positive Vorhersagen macht.
	\item \textbf{recall (Trefferquote)}: Der Anteil der korrekt vorhergesagten positiven Instanzen an allen tatsächlichen positiven Instanzen. Ein hoher Recall bedeutet, dass das Modell wenige falsch negative Vorhersagen macht.
	\item \textbf{f1-score}: Das harmonische Mittel aus Präzision und Recall. Es gibt einen einzigen Wert an, der die Balance zwischen Präzision und Recall berücksichtigt.
	\item \textbf{support}: Die Anzahl der tatsächlichen Vorkommen der Klasse im Testset. Es gibt an, wie viele Instanzen jeder Klasse im Testset vorhanden sind.
	\item \textbf{accuracy (Genauigkeit)}: Der Anteil der korrekt vorhergesagten Instanzen an allen Instanzen. Es gibt an, wie oft das Modell insgesamt richtig liegt.
	\item \textbf{macro avg}: Der Durchschnitt der Metriken (Präzision, Recall, f1-score) über alle Klassen hinweg, ohne Berücksichtigung der Klassenhäufigkeit.
	\item \textbf{weighted avg}: Der gewichtete Durchschnitt der Metriken (Präzision, Recall, f1-score) über alle Klassen hinweg, wobei die Klassenhäufigkeit berücksichtigt wird.
\end{itemize}

\subsection{Diskussion der Ergebnisse}
Die in Abbildung ref{tab:klassifikationsreport} dargestellten Ergebnisse des Klassifikationsreports zeigen, dass das Modell besonders gut dazu in der Lage ist, die Aktivitäten "Gehen\dq\, und "Joggen\dq\, zu klassifizieren. Die Aktivität "Rennen\dq\, wurde ebenfalls gut erkannt, allerdings etwas schlechter als die Aktivitäten "Gehen\dq\, und "Joggen\dq. Die hohe Genauigkeit (Accuracy) von 0,95 sowie die hohen Werte für den makro- und gewichteten Durchschnitt (Macro Avg und Weighted Avg) von jeweils 0,95 bzw. 0,94 zeigen, dass das Modell insgesamt eine sehr gute Leistung erbringt. Für zukünftige Arbeiten könnte man sich darauf konzentrieren, noch mehr Daten von verschiedenen Testpersonen zur Aktivität "Rennen\dq\, zu sammeln, um die Leistung des Modells zu verbessern. Unsere Datenmenge sollte jedoch ausreichen, um sagen zu können, dass es sich um gute Ergebnisse handelt, da „Rennen” wahrscheinlich als schnelles "Joggen\dq\, klassifiziert wird und somit schwer zu unterscheiden ist. %Quelle hier hinzufügen

\bgroup
\def\arraystretch{1.5}
	\begin{table}[h]
		\centering
		\begin{tabular}{|l|c|c|c|c|}
			\hline
			& precision & recall & f1-score & support \\
			\hline
			gehen & 0.96 & 0.99 & 0.97 & 148 \\
			\hline
			joggen & 0.94 & 0.93 & 0.93 & 98 \\
			\hline
			rennen & 0.94 & 0.91 & 0.92 & 98 \\
			\hline
			accuracy & & & 0.95 & 344 \\
			\hline
			macro avg & 0.95 & 0.94 & 0.94 & 344 \\
			\hline
			weighted avg & 0.95 & 0.95 & 0.95 & 344 \\
			\hline
		\end{tabular}
		\caption{Klassifikationsreport}
		\label{tab:klassifikationsreport}
	\end{table}
\egroup

\subsection{Spätere Wiederverwendung des Modells}
Die trainierten Modelle werden anschließend mithilfe der joblib-Bibliothek gespeichert, um sie später wieder verwenden zu können. Dies ermöglicht eine effiziente Wiederverwendung der Modelle ohne erneutes Training. 

